\chapter{Development Plan \\
\small{\textit{- Ryan Davis, Cory Vitanza, and Roy Tas}}}
\index{development plan} 
\index{Chapter!Development Plan}
\label{Chapter::DevelopmentPlan}

\section{Introduction}

The purpose of this project is to design, implement, and validate a location-tracking and data communication system using LoRaWAN-enabled hardware integrated with a cloud-based software platform. The system aims to provide reliable long-range wireless data transmission between embedded devices and a backend service, with supporting web-based tools for monitoring and management. By combining embedded development, cloud infrastructure, and a user-facing interface, this project demonstrates the end-to-end development process of a distributed system.  

Success criteria for the project include (1) establishing reliable LoRaWAN communication between the embedded board and backend, (2) ensuring data is properly stored, processed, and displayed through the frontend application, and (3) passing acceptance testing with the Wearable Robotics Systems Lab. From a user perspective, the system enables simple, reliable collection and visualization of data in real time, serving as a proof-of-concept platform that can be extended to future applications.


\section{Roles and Responsibilities}

Clear definition of roles and responsibilities ensures accountability and effective collaboration within the team.  
Since this project is being developed by a small group, several members will serve multiple roles across development, testing, documentation, and risk management.  
The following assignments identify primary responsibilities, while recognizing that all team members will contribute collaboratively to design decisions, implementation, and review processes.

 
\begin{enumerate}
\item Development Lead: Ryan
\item Buildmeister: Roy
\item Architect: Cory
\item Developers: Ryan (Backend), Cory (Frontend), Roy (Infrastructure)
\item Test Lead: Cory
\item Testers: All
\item Documentation: All
\item Documentation Editor: All
\item Designer: Cory
\item User advocate: Ryan Davis
\item Risk Management: Roy
\item System Administrator: Ryan
\item Modification Request Board: All
\item Requirements Resource: Roy
\item Customer Representative: All
\item Customer responsible for acceptance testing: Wearable Robotics Systems Lab
\end{enumerate}


\section{Method}
These are unique to software development, although there may be some overlap.

\subsection{Software}
\begin{enumerate}
\item Language(s): 
  \begin{itemize}
    \item C (Arduino IDE)
    \item Python (3.x)
    \item TypeScript (React)
    \item SQL (PostgreSQL via Supabase)
  \end{itemize}

\item Operating System(s):
  \begin{itemize}
    \item Arduino embedded OS (board-specific)
    \item Linux/Ubuntu (for Docker and AWS deployment)
    \item macOS/Windows (development environments)
  \end{itemize}

\item Software packages/libraries used:
  \begin{itemize}
    \item LoRaWAN protocol (TCP/IP stack integration)
    \item Supabase (PostgreSQL backend service)
    \item Google Maps API
    \item React (TypeScript frontend framework)
    \item Docker (containerization)
    \item AWS (deployment and hosting)
    \item GitHub (version control)
    \item Kanban Board (Agile task management)
  \end{itemize}

\item Code conventions: 
  \begin{itemize}
    \item GitHub repository guidelines for commits and pull requests
    \item ESLint + Prettier for TypeScript (frontend)
    \item PEP 8 for Python (backend tools)
    \item Arduino C style conventions
  \end{itemize}
\end{enumerate}


\subsection{Hardware}
\begin{enumerate}
\item Development Hardware 
    \begin{itemize}
        \item Wio Tracker 1110 Development Board
        \item Semtech LR1110 LoRa® transceiver
    \end{itemize}

\item Test Hardware 
    \begin{itemize}
        \item Wio Tracker 1110 Development Board
        \item LoRaWAN Gateway (connected to The Things Stack)
    \end{itemize}

\item Target/Deployment Hardware 
    \begin{itemize}
        \item Wio Tracker 1110 Development Board
        \item Semtech LR1110 LoRa® transceiver
        \item LoRaWAN Gateway (The Things Stack integration for live deployment)
    \end{itemize}
\end{enumerate}


\subsection{Backup plan (individual and project)} 

In case the Wio Tracker is not capable of the deployment of our software, the FindIt team will need to research and develop board alternatives. One option for this scenario would be to use a Raspberry Pi Nano, implemented with an external LoRaWAN module to allow for data transmission. Even with such addition, it is necessary for our unit to maintain battery consumption amongst boards, ensuring our unit has extensive battery life for users.

\subsection{Review Process} 
\begin{enumerate}
\item Will you do architecture, usability, design, security, privacy or code reviews?
    \begin{itemize}
        \item Code reviews will be completed upon the completion of Agile tasks on our Kanban board.
    \end{itemize}

\item What approach will you use for the reviews (formal, informal, corporate standard)? 
    \begin{itemize}
        \item Informal reviews will be employed to achieve the approval of all members of the team, ensuring that all cases are handled in new implementation.
    \end{itemize}
\item Who is responsible for the reviews and resolving any issues uncovered by the reviews? 
    \begin{itemize}
        \item Upon the implementation of a component, the other two members of the team will perform a code review. Any issue uncovered from code review will be resolved by the member working on the component.
    \end{itemize}
\end{enumerate}

\subsection{Build Plan} 
\begin{enumerate}
\item Revision control system and repository used
    \begin{itemize}
        \item GitHub (Version Control)
    \end{itemize}
\item Continuous integration 
    \begin{itemize}
        \item Circle CI
    \end{itemize}
\item Deadlines for the builds
    \begin{itemize}
        \item Depending on our planned sprints, our Kanban board will present deadlines based on tasks.
    \end{itemize}
\item Regression test process – see test plan 
\end{enumerate}

\subsection{Modification Request Process} 
\begin{enumerate}
\item MR tool: For modification requests, we will use our Kanban board on GitHub to issue any modification requests as a task in the board.
\item Decision process: Decisions for modification will be decided during heartbeat meetings, after a discussion among the team concerning the modification. The team will proceed to assign a task to the Kanban board on the methodology and developers necessary for the task.
\item State whether there will be two process streams one during development and one after development:
    \begin{itemize}
        \item During development, there will be two process streams. One is this programming of the development board and the other will be for the front-end user interface. The two streams will be integrated for deployment.
    \end{itemize}
\end{enumerate}

\section{Virtual and Real Workspace}
\subsection{Virtual Workspace}
Our team will utilize texts for direct communication, as well as a Discord server used for the sharing of files, resources, and communication. GitHub will also be used for secondary communication, implementing comments and changes to the code base via commits and merge requests.
\subsection{Real Workspace}
The team's regular workspace will be the Software Engineering Lab, primarily for heartbeat meetings and the development of our program. Meetings with Dr. Zanotto will be held in his office or over Zoom, when necessary.

\section{COMMUNICATION PLAN}

\subsection{Heartbeat Meetings}
Heartbeat meetings will be held at least twice per week with all group members present.  
These meetings serve to take the “pulse” of the project by providing quick updates, raising blockers, and reviewing open issues and risks.  
Each meeting will follow a structured agenda that includes:  
\begin{itemize}
    \item Individual progress updates  
    \item Discussion of current blockers or risks  
    \item Review of action items and open issues  
\end{itemize}
Notes from each meeting will be recorded in the project document, and the GitHub Kanban board will be updated accordingly.  
These meetings are designed to be short, ideally during our allotted class-time.

\subsection{Status Meetings}
Status meetings will occur once a month and will include the project group and Dr.~Damiano Zanotto.  
The purpose of these meetings is to present overall project progress, demonstrate completed work, and receive feedback on areas for improvement.  
If issues arise during these meetings, they will be addressed separately in an issues meeting.  
These meetings are intended to be concise, focusing on high-level updates and progress checkpoints.

\subsection{Issues Meetings}
Issues meetings will be scheduled whenever significant problems arise that require group consensus or immediate attention.  
Alerts for these meetings will typically be raised through a Discord message.  
If consensus is reached that a meeting is necessary, it will be scheduled as soon as all group members are available.  
These meetings provide a dedicated space for problem-solving and decision-making outside of the regular heartbeat meetings.


\section{Timeline and Milestones}
The following milestones outline the major deliverables for the project. Each milestone 
represents a 100\% complete item, with critical participants, begin and end dates, and 
deliverables clearly defined.

\begin{itemize}
    \item \textbf{Milestone 1: Project Setup \& Requirements Finalization}  
    \begin{itemize}
        \item Begin: Oct 3, 2025 \quad End: Oct 17, 2025  
        \item Deliverables: Project repository initialized, Kanban board configured, requirements document finalized, hardware/software resources received and confirmed.  
    \end{itemize}

    \item \textbf{Milestone 2: Core System Architecture \& Prototype Communication}  
    \begin{itemize}
        \item Begin: Oct 17, 2025 \quad End: Oct 31, 2025  
        \item Deliverables: Draft system architecture, initial Arduino board programming for basic LoRaWAN communication, mock frontend interface, database schema in Supabase.  
    \end{itemize}

    \item \textbf{Milestone 3: Integration of Hardware and Software Components}  
    \begin{itemize}
        \item Begin: Oct 31, 2025 \quad End: Nov 14, 2025  
        \item Deliverables: Hardware successfully transmitting data to backend, frontend connected to database, Docker setup for local deployment, integration testing initiated.  
    \end{itemize}

    \item \textbf{Milestone 4: Feature Implementation \& System Testing}  
    \begin{itemize}
        \item Begin: Nov 14, 2025 \quad End: Nov 28, 2025  
        \item Deliverables: Full set of planned features implemented, CI/CD pipeline operational, unit + integration tests passing, preliminary user documentation drafted.  
    \end{itemize}

    \item \textbf{Milestone 5: Final Review \& Delivery}  
    \begin{itemize}
        \item Begin: Nov 28, 2025 \quad End: Dec 12, 2025  
        \item Deliverables: Final acceptance testing completed, documentation finalized, project demonstration delivered, final repository and deployment package prepared for submission.  
    \end{itemize}
\end{itemize}

\noindent
This document will be re-issued at each milestone to reflect progress and highlight 
changes since the previous revision.


\section{Testing Policy/Plan}
Our testing policy combines automated and manual approaches to ensure that the project 
meets functional requirements and maintains high reliability. The following types of testing 
will be performed:

\begin{itemize}
    \item \textbf{Unit Testing:} We will implement unit tests for critical components, 
    including geofence logic and location-based features. Pytest will be the primary 
    testing framework. Unit testing will begin as soon as core modules are developed 
    and will continue throughout the project.
    
    \item \textbf{Integration Testing:} Integration testing will verify interactions between 
    different system components, such as API calls to The Things Stack and web service 
    calls to the Google Maps API. These tests will be mocked where appropriate to ensure 
    reliability and repeatability.

    \item \textbf{System Testing:} End-to-end system testing will be conducted once a 
    minimum viable product is available. This will include route pings, location checking, 
    and verifying that hosted meals can be created, displayed, and joined.

    \item \textbf{Acceptance Testing:} Acceptance testing will be performed in collaboration 
    with Dr.~Damiano Zanotto. This testing will validate that the system meets the agreed 
    project requirements and is ready for delivery.

    \item \textbf{Regression Testing:} Regression tests will be rerun whenever bug fixes or 
    new features are introduced to ensure that existing functionality remains stable. 

    \item \textbf{Testing for Quality Attributes:} While performance and security are not the 
    primary focus, geolocation accuracy and reliability will be emphasized. Specific attention 
    will be given to ensuring that geofences trigger consistently and API interactions are 
    fault-tolerant.

    \item \textbf{Continuous Integration (CI):} CircleCI will be used to automate testing during 
    development. All code changes will trigger automated test runs to catch issues early.
\end{itemize}

\noindent
Our stopping criteria for testing is defined as achieving at least 80\% unit test coverage across critical modules, successful integration of geofencing and API functionality, and passing acceptance testing with project stakeholders.


\section{Risks}
Risks that we must mitigate throughout the development of our product include various issues on all scales of development. With the development board, we want to ensure that the battery life is competitive with existing tracking units, such as AirTag. We also would like to ensure that the tracking unit is able to provide precise location data for the user to accurately locate the unit.

\section{Assumptions (required)}
For the development of our program, it is assumed that:
\begin{itemize}
    \item Publicly accessible LoRaWAN gateways populate our region, allowing for fluid connection between our board to LoRaWAN network servers.
    \item End users have mobile devices or computers that are able to access our website, providing them access to their tracking units.
\end{itemize}

\section{DISTRIBUTION LIST}
 This document will be shared with the members of the FindIt team, as well as Dr. David Darian Muresan and Dr. Damiano Zanotto throughout our development processes.

\section{IRB Protocol}
Our project does not require IRB protocol.


\section{Required Resource and Budget}

\subsection{Hardware Resources}
For our tracking unit, the only development expense will include the Wio Tracker 1110 development board, costing \$30.
\subsection{Software Resources}
Depending on the size of our infrastructure, deployment on AWS as well as the storage of historical data on Supabase will require subscriptions in order to maintain the usage of our software.

